\section{Symbols and Theology under the Rule of the Gospel}

\subsection{The crucifix is the source of light}

The most important visual claim is not the mountain, clothing, tears, roses, or instruments. The children say the light radiated from the crucifix on the Lady's breast. The reported apparition is therefore iconographically Christocentric before any interpretation begins. Mary is visible in light received from the crucified Christ. Her speech points to ``my Son''; her mission is intelligible only within his Paschal mystery.

This accords with Catholic doctrine. Christ is the one Mediator and Redeemer (1 Tim 2:5--6). Mary's cooperation is created, graced, maternal, and wholly dependent on him (LG 56, 60--62; CCC 970). She neither adds saving merit to an insufficient Cross nor controls access to an unwilling Savior. The DDF's 2025 \emph{Mater Populi Fidelis} insists that Marian language must preserve this order. Any reading of La Salette that leaves Christ at the edge while turning Mary into an autonomous ruler has inverted the reported sign.

\subsection{The Son's arm and the mercy of the Judge}

The ``heavy arm'' is La Salette's most theologically dangerous image when detached from Scripture. Biblical ``arm'' language can signify divine power to save, govern, and judge: the Lord brings Israel out with an outstretched arm; Isaiah sings of the Lord's holy arm; Mary's Magnificat praises the strength of God's arm. It is not anatomical information about the risen Christ.

The reported sentence combines judgment and intercession. Persistent contempt for God and neighbor has consequences; grace does not declare evil unreal. Yet Mary cannot be cast as the compassionate party restraining a cruel Son. The Son whose judgment is serious is the same Son who became flesh, wept, forgave enemies, bore sin, and gave himself for the world. His ``arm'' is not only the power before which sin is exposed; it is the strength by which the sinner is rescued.

John Paul II's 1996 letter for the 150th anniversary reads the message as hope. He interprets Christ's arm through mercy: it does not condemn those who walk humbly and reaches toward the sinner. He describes Mary's prayer as continual intercession that brings her children toward the Son. This papal reception does not rewrite the historical words; it establishes the Catholic theological horizon in which they may safely be preached.

\interpretivekey{No rivalry within divine mercy}{Mary does not persuade Christ to become merciful. Her intercession is itself a fruit and instrument of the triune God's mercy. The Son freely wills his Mother's prayer and the communion of saints; she freely participates in his one mediation. ``Holding the arm'' can express the urgency of maternal intercession, never a contest between a tender Mother and a vindictive Redeemer.}

\subsection{The tears}

The figure is seated, face in hands, and reportedly weeps throughout the discourse. Tears place warning inside lament. Scripture permits this grammar: prophets mourn the people; Jesus weeps over Jerusalem and at Lazarus's tomb; Paul writes with tears. Sorrow does not deny hope. It reveals that sin wounds communion and that warning can be an act of love.

Devotion may call Mary the Weeping Mother or recall her compassion at the Cross. It should not construct a physiology of heavenly grief or imply that the glorified Mother of God lacks beatitude. The vision is symbolic mediation. Its tears translate maternal compassion into a form the children and Church can contemplate. Nor are they emotional leverage: fear or guilt that isolates a person from Christ is not their proper fruit. Their authentic fruit is compunction joined to hope and mercy toward other sufferers.

\subsection{Hammer, pincers, chain, and roses}

The children described a hammer and pincers beside the corpus. They did not report the Lady explaining them. Later iconography often reads the hammer as sin driving the nails and the pincers as repentance, reparation, or charity removing them. That is a coherent devotional contrast, provided it is labeled as interpretation rather than reported speech. No creature literally repeats Calvary, and Christian reparation does not complete a deficient sacrifice. The faithful allow Christ's once-for-all gift to bear fruit by turning from what wounds his members and by repairing harms they have caused.

The heavy chain and garlands of roses admit several spiritual readings. A chain can suggest bondage to sin, solidarity with labor, or the weight of suffering; roses can suggest grace, beauty, prayer, and the life that surrounds the Cross. The local dress makes the figure near to Alpine working people. None of these meanings was verbally assigned in the early public message. Their theological value depends on conformity to Scripture and the Church, not imaginative certainty.

\begin{threestable}{Reported sign}{Responsible theological reading}{Overclaim to reject}
Light from the crucifix & Christ's Cross and Resurrection are the source and center; Mary shines by received grace. & Mary as an independent source of salvation or revelation.\\
Tears & Maternal compassion, prophetic lament, and a call to compunction. & Manipulative terror or a claim that heavenly beatitude is defective.\\
Arm & Divine power, the seriousness of sin, judgment ordered to conversion, and saving strength. & Mary protecting humanity from a Son or Father less merciful than she.\\
Hammer and pincers & Later devotional contrast between sin and grace-enabled repentance or repair. & A divinely dictated symbol-key or human works completing Calvary.\\
Chain and roses & Bounded typology of bondage, solidarity, beauty, prayer, and grace around the Cross. & One mandatory decoding or proof of a later devotion not named in 1846.\\
Mountain and ascent & Biblical place of encounter and an image of transcendence and mission. & Geography as supernatural proof or pilgrimage as a condition of salvation.\\
\end{threestable}

\subsection{Sunday: gift before obligation}

The reported complaint about six days and the seventh echoes the Decalogue (Exod 20:8--11; Deut 5:12--15). Because the speaker is Mary, ``I gave you six days'' must be read as prophetic representation of God's claim, not a declaration that Mary created time or promulgates divine law on her own authority. Biblical messengers frequently voice God's word in the first person. The Church's law and doctrine, not a private revelation, establish the obligation.

For Christians the Lord's Day is the weekly Pasch: the day of the Resurrection, the Eucharistic assembly, worship, joy, rest, family, and works of mercy. La Salette's complaint is not merely about a breached attendance rule. It exposes a life in which production crowds out worship, the sacred Name becomes verbal debris, and Mass becomes spectacle. Conversion restores relationship.

The warning cannot be used against those compelled by poverty, care, essential services, illness, disability, persecution, or absence of accessible ministry. Catholic teaching itself recognizes necessary family and social service. Employers, consumers, and communities share responsibility for conditions that prevent rest. A La Salette spirituality of Sunday asks not only ``Did you work?'' but also ``Whose labor makes my leisure possible, and do the poor receive rest?''

\subsection{The holy Name, Mass, and Lent}

The second named offense is profaning Christ's name. This belongs to the command not to take the Lord's name in vain and to the New Testament discipline of truthful speech (Mt 5:33--37; Jas 5:12). The message names cart drivers because their oaths were concrete and audible, not because one occupation is uniquely sinful. Modern application includes casual blasphemy, religious speech used to threaten, false oaths, and invocation of God to sanctify hatred or deceit.

The complaints about Mass and Lent reveal practical unbelief. Going to church merely to mock, or treating penitential discipline with contempt, disconnects culture from conversion. Yet private revelation cannot freeze 1846 French discipline. The faithful follow the Church's present calendar, precepts, fasting and abstinence law, legitimate dispensations, and pastoral care. A sick person, child, pregnant person, laborer, or person with an eating disorder must never be coerced into harmful fasting by appeal to La Salette.

\subsection{Famine, judgment, and the innocent sufferer}

The crop and famine passages can be read irresponsibly in two opposite ways. One can remove moral agency entirely and hear only primitive meteorology. Or one can assert that every blight, illness, or dead child is a direct penalty for identifiable personal sins. Scripture permits neither simplification.

Biblical covenant and prophetic texts connect communal injustice and idolatry with the unraveling of life. Human conduct truly can worsen hunger: hoarding, exploitation, war, neglect, unjust land arrangements, disregard for creation, and indifference to the poor have material effects. A community's spiritual disorder can become social and ecological disorder. La Salette's agricultural language preserves that moral seriousness.

Jesus, however, refuses the inference that particular sufferers are necessarily worse sinners (Lk 13:1--5), and he rejects the attempt to assign a man's blindness to his or his parents' sin (Jn 9:1--3). Job dismantles easy retribution. The child trembling in a caregiver's arms is first a neighbor to protect, feed, mourn, and love---never an exhibit in someone else's moral argument. The reported tears stand over the famine warning and forbid cold blame.

The message's ``if converted'' turn also places prophecy in the biblical field of conditional appeal (Jer 18:7--10; Jonah). Its aim is not to satisfy curiosity about a fixed future but to change the present. Failure to map every image to a later statistic is not failure of faith. Conversion, providence, weather, human freedom, and economic systems cannot be reduced to a private-revelation equation.

\subsection{Prayer small enough to begin}

When asked whether they pray, the chosen witnesses admit that they scarcely do. The prescribed minimum---an Our Father and Hail Mary morning and evening when they cannot do better---is pastorally gentle. The message does not require immediate spiritual virtuosity. It gives a beginning that a child can remember.

The Our Father orders desire to God's Name, Kingdom, will, daily bread, forgiveness, and deliverance from evil. In the context of La Salette, ``daily bread'' is not metaphor only. The Hail Mary joins Scripture's greeting to a request for intercession now and at death. More prayer is encouraged when possible, but quantity is not a currency that purchases immunity from suffering (CCC 2111). Prayer opens the person to grace, truth, endurance, and love.

\subsection{Reconciliation as the mature synthesis}

The title ``Reconciliatrix'' was not spoken in the encounter, but it became a recognized devotional synthesis. The Holy Office's 1915 decree expressly preserved devotion to the Blessed Virgin under the title commonly called Reconciliatrix of La Salette even while prohibiting discussion of the secret. The Missionaries define their charism as service of Christ and the Church in the mystery of reconciliation.

The title must be governed by 2 Corinthians 5:18--20 and Colossians 1:19--20: God reconciles the world to himself in Christ and entrusts the ministry of reconciliation to the Church. Mary participates as the Mother whose \latin{fiat} receives the Redeemer, whose intercession is maternal and subordinate, and whose reported message sends sinners back to her Son. She is not a second author of the Atonement.

Reconciliation is therefore more demanding than sentiment. It includes confession and absolution where appropriate, restitution, truth-telling, repair of scandal, forgiveness that does not deny justice, food for the hungry, protection of children, Sabbath rest for workers, and reunion with the Eucharistic body of Christ. It never requires a victim to return to danger, suppress evidence, forgo safeguarding, or call abuse peace. The Cross exposes evil as it forgives sinners.

\interpretivekey{The public message has no apocalypse to decode}{Its urgency is already complete without the later secret literature: worship God, reverence Christ's Name, pray, repent, honor Sunday, take hunger seriously, and communicate hope. The public message gives conditions and practices, not a calendar of popes, nations, Antichrist, or the end of the world. Those themes belong to separately dated texts and must be judged there.}
